% Outline preserved for next pass:
%
% Five mandatory bullets this abstract must hit:
%   1. Problem: heterogeneous eval-suite mixes scorer classes × task shapes;
%      single-axis framing collapses three distinct reliability axes.
%   2. Current-practice limitation: dominant open-source eval stacks default
%      to RLHF-tuned instruction-following judges — precisely the configuration
%      the calibration literature documents as miscalibrated.
%   3. Approach (a): three-axis methodology — calibration-no-regression,
%      agreement-no-regression, bias-bound precondition.
%   3. Approach (b): four surgical v1 upgrades — probability-weighted scoring,
%      position-bias swap test, ECE + temperature scaling, stratified sampling.
%   4. Headline result: preliminary empirical results on a held-out calibration
%      set; no fabricated numbers.
%   5. Contribution: three-axis methodology + per-cell priors matrix +
%      architecture + v1 algorithm + eval pipeline.

\begin{abstract}
\noindent
Multi-agent coding frameworks expose a heterogeneous eval-suite: five scorer
classes (heuristic, embedding, LLM classifier, LLM rubric, trajectory)
applied across eight task shapes (single-file edits, Q\&A lookups, ADR
judgments, threat models, long-form docs, code review, debug triage, rubric
writing).
The dominant single-axis framing in the eval literature --- either
judge-to-human agreement or expected calibration error --- collapses three
distinct reliability axes that fail in different cells of the scorer-class
$\times$ task-shape grid for structurally different reasons and require
different fixes.

A compounding operational problem is that the eval tools most commonly
deployed --- forced-tool-use LLM-as-judge on RLHF-tuned instruction-following
models --- occupy precisely the judge configuration that the calibration
literature already documents as systematically miscalibrated relative to
pretrained models of the same scale.

We propose a three-axis methodology: (i)~\emph{calibration-no-regression},
for scorers that emit a confidence-vs-correctness signal (LLM classifier,
pairwise rubric); (ii)~\emph{agreement-no-regression}, for scorers whose only
stable signal is inter-rater consistency (LLM rubric, trajectory rubric on
subjective task shapes); and (iii)~a \emph{bias-bound precondition} on every
LLM-judge cell --- position-bias flip-rate $\leq 5\%$, verbosity-bias
correlation $|\rho| \leq 0.20$, self-preference advantage $\leq 0.05$.
We anchor per-cell axis tolerances to published ML evaluation priors.
Atop this frame we ship four surgical v1 upgrades: probability-weighted
scoring (G-Eval; Liu et al.\ EMNLP 2023) recovering calibration signal
already present in judge log-probs, a position-bias swap test
(Zheng et al.\ NeurIPS 2023), ECE diagnostics with temperature scaling
(Guo et al.\ ICML 2017), and stratified sampling for production-trace suites.

We report preliminary results on a held-out labelled calibration set per
judge model; full empirical results will follow as paper-stream telemetry
accumulates from the Blackrim judge fleet.
Probability-weighted scoring is the headline implementation finding: a
$\sim$30-line change that, in the literature, recovers a 22\% Spearman
correlation gain over argmax decoding at no additional API cost.

The core contribution is a per-cell reliability methodology that treats
calibration, agreement, and bias bounds as non-collapsible axes, paired with
an eval-suite architecture and paper-stream telemetry pipeline that makes the
three-axis gate operationally tractable as a release-blocking CI check.
\end{abstract}
